\scanpage{367}\mspage{283}
D'ailleurs, les conditions \unsure{($\alpha$)|(a)} impliquent que
$\pi\deleted{\unsure{?}}$ est utile pour $\underline S$, donc
\[
 \underline S\xleftarrow{\approx}\mathrm{Ind}_\pi(S_\pi)\unsure{,|;}
\]

% margin note written sideways
\margin{``facultatif''}
[D'ailleurs, \deleted{\uuline{\unsure{comme}}}\inserted{si} \uline{$\pi'_0\geq\pi_0$}, la relation
$\pi^{\pi'_0}=\pi$ implique a fortiori $\pi^{\pi_0}=\pi$,
et \deleted{\unsure{bien sûr}}\inserted{si} \uline{$\pi'_1\geq\pi_1$}, on voit que
($\alpha$) implique que $\pi$ est utile aussi
pour $M$, i.e.
\[
 M\xleftarrow{\approx}\mathrm{Ind}_\pi(M_\pi).
\]
D'autre part, le foncteur d'inclusion
$\underline S\subset M$ se déduit \inserted{\unsure{essentielt|essentiellement}} de $\underline S_\pi\subset M_\pi$
en passant aux catégories $\mathrm{Ind}_\pi$.]

\margin{$\unsure{\surd}$} Ceci prouve \inserted{entre autres} le ``il faut'' dans la
proposition qui suit\unsure{:}

\uline{Proposition} 4.13.2. Soient $M$ une
cat. \deleted{\unsure{accessible}}\inserted{accessible}, $\pi_0$ un cardinal infini
\deleted{supposons}\inserted{utile pour} $M$, \deleted{$\pi_0$-accessible}. $\unsure{S|\mathscr S}$ une
\deleted{\unsure{sous-caté}} partie de $\operatorname{Ob}M$, \deleted{Pour que}\inserted{$\underline S$ la}
\deleted{$S$ soit accessible}\inserted{sous-catégorie \unsure{pleine} de $M$ correspondante},
% archive sheet number: 366

